\documentclass[12pt]{article}
\usepackage[T1]{fontenc}
\usepackage[utf8]{inputenc}
\usepackage{lmodern}
\usepackage{textcomp}
\usepackage{enumitem}
\usepackage{geometry}
\geometry{margin=1in}
\begin{document}
\begin{center}
Answer Key - Philosophy - Undergraduate Level Assessment 15 \
\small (Answers are ordered exactly as in the question paper.)
\end{center}

\section*{Multiple Choice}
\begin{enumerate}[label=\arabic*.]
\item \textbf{Answer:} A
\\ \textit{Options:} A) the artificial assumptions are, in fact, empirical claims that are central to the debate.; B) it leads us to believe that torture is always permissible.; C) all artificial thought experiments are philosophically useless.
\\ \textit{Note:} Marcia Baron highlights that the artificial assumptions in ticking bomb cases are not merely theoretical constructs but rather empirical claims that fundamentally influence the ethical debate surrounding torture and its justification.
\item \textbf{Answer:} A
\\ \textit{Options:} A) anarchy.; B) deontological theory.; C) libertarianism.; D) communitarianism.; E) socialism.
\\ \textit{Note:} Wellman argues that insisting on unlimited property rights leads to anarchy because it undermines the social cooperation necessary for a functioning society, as individuals would prioritize their own interests over communal order, resulting in chaos and conflict.
\item \textbf{Answer:} A
\\ \textit{Options:} A) Mengzi; B) Zoroastrianism; C) Zen Buddhism; D) Shintoism; E) Vedic Philosophy
\\ \textit{Note:} Mengzi, also known as Mencius, emphasized the inherent goodness of human nature and the importance of aligning with natural processes, which aligns with the principles of naturalism in philosophy.
\item \textbf{Answer:} A
\\ \textit{Options:} A) worldly possessions and wealth; B) acts of the moral virtues.; C) loving God.; D) pleasure; E) achieving personal success
\\ \textit{Note:} The correct answer is based on Aquinas's belief that while ultimate happiness is found in the divine and eternal, he acknowledges that many people mistakenly equate happiness with worldly possessions and wealth, reflecting a common misunderstanding of true fulfillment.
\item \textbf{Answer:} A
\\ \textit{Options:} A) Zhuangzi; B) Dao; C) Xunzi; D) Laozi; E) Zisi
\\ \textit{Note:} Zhuangzi is considered a mystic because his philosophical writings emphasize the concept of qi, exploring the nature of reality and the interconnectedness of all things through a lens of spontaneity and transformation.
\end{enumerate}

\section*{True / False}
\begin{enumerate}[label=\arabic*.]
\setcounter{enumi}{5}
\item \textbf{Answer:} True
\\ \textit{Options:} A) True; B) False
\\ \textit{Note:} According to Kant, the laws of nature pertain to the deterministic processes of the physical world, while laws of freedom relate to moral imperatives that guide human actions based on rationality and ethical considerations, thus making the statement true.
\item \textbf{Answer:} False
\\ \textit{Options:} A) True; B) False
\\ \textit{Note:} The 'Flower Sermon' is primarily associated with Zen Buddhism, where it emphasizes direct transmission of insight rather than a focus on textual recitation, distinguishing it from Nichiren Buddhism, which centers on the Lotus Sutra.
\item \textbf{Answer:} False
\\ \textit{Options:} A) True; B) False
\\ \textit{Note:} Singer argues that the loss of human diversity is a significant concern, but he emphasizes that the ethical implications of increased disparity between the rich and the poor are equally, if not more, serious due to the profound effects on human well-being and social justice.
\item \textbf{Answer:} False
\\ \textit{Options:} A) True; B) False
\\ \textit{Note:} The statement is false because Smith's action involves an active decision to not intervene, which contrasts with the concept of letting die, where one passively allows an outcome without taking any action to influence it.
\item \textbf{Answer:} True
\\ \textit{Options:} A) True; B) False
\\ \textit{Note:} Mill argues that utilitarianism can conflict with our intuitive sense of justice, particularly when actions that maximize overall happiness may violate individual rights or lead to unjust outcomes, highlighting a significant tension between utilitarian principles and notions of justice.
\end{enumerate}

\section*{Long Form Response}
\begin{enumerate}[label=\arabic*.]
\setcounter{enumi}{10}
\item \textbf{Answer:} The externalist approach to meaning aligns with existentialist philosophy by emphasizing the significance of context, relationships, and the external world in shaping individual meaning and experience. Existentialists, like Jean-Paul Sartre and Simone de Beauvoir, argue that existence precedes essence, suggesting that individuals must create their own meanings through interactions with their environment and others. This perspective underscores the idea that human experience is not solely an internal, subjective affair but is deeply intertwined with external factors, such as societal norms and interpersonal relationships. Consequently, understanding meaning as external invites a broader exploration of how individuals navigate their existence, highlighting the dynamic interplay between personal agency and the surrounding world in defining what it means to be human.
\\ \textit{Note:} correct because it connects the externalist approach to meaning with existentialist philosophy by highlighting how context, relationships, and the external world shape individual meaning, reflecting the existentialist idea that individuals must create their own meanings through their interactions and experiences.
\item \textbf{Answer:} Van den Haag argues that murderers who receive the death penalty are morally degraded due to their heinous actions, which reflect a profound disregard for the sanctity of human life. This degradation raises questions about whether we can truly ascertain their moral status, as their actions may stem from various factors, including psychological issues, societal influences, or personal circumstances. The implication for justice is significant; if we accept that moral degradation exists, we must consider whether the death penalty serves as a just response or merely perpetuates a cycle of violence. Ultimately, the challenge lies in balancing the need for justice with the recognition of the complex moral landscapes that define human behavior.
\\ \textit{Note:} The answer correctly identifies Van den Haag's view that murderers display moral degradation through their actions, while also acknowledging the complexities in determining their moral status, which raises important implications for the justice system's approach to punishment.
\end{enumerate}

\vfill
\noindent\textit{End of Answer Key}
\end{document}